% ============================================================
%  文档类选项（在 main.tex 的 \documentclass[...]{tongjithesis} 中设置）
%  field=science    理工科（工科类、理科类，默认）— 章节编号：1 / 1.1 / 1.1.1
%  field=humanities 文科  （人文、法学、外语、艺术）        — 章节编号：一、/（一）/ 1.
%  经管类属社科：编号要求依理工，撰写依理工模版；26 届为过渡期，可选 science 或
%  humanities；自 27 届起统一使用 science。
% ============================================================

% ============================================================
%  封面信息（请根据实际情况填写）
% ============================================================
\school{计算机科学与技术学院}
\major{信息安全}
\student{2252552}{胡译文}
\thesistitle{基于状态机的自动化渗透测试框架设计与实现}{}
\thesistitleeng{Design and Implementation of Automated Penetration Testing System Based on Finite State Machine}{}
\thesisadvisor{钟计东}
\thesisdate{\the\year}{\the\month}{\the\day}

% ============================================================
%  信息说明页
% ============================================================

% 成果类型（三选一）：
%   thesis      — 毕业论文，摘要
%   design      — 毕业设计（软件/创意/建筑等非工程设计类），摘要
%   engineering — 毕业设计（工程设计类），设计总说明
\infotype{thesis}

% 内容简述（请用300字以内简要概述）
\infoabstract{
本文针对基于大语言模型的自动化渗透测试框架AutoPT在实际应用中的不足，从流程控制、工具能力和交互方式三个方面进行了改进。在流程控制方面，设计了统一的ReAct结构化提示词模板与角色差异化机制，将Scan阶段细化为四个子阶段并引入提前终止机制，实现了上下文压缩和PoC容错提取，改进Check阶段的验证逻辑。在工具能力方面，新增三个自定义浏览器扩展工具实现了表单交互能力，实现了多级容错的ReAct解析器。在交互方式方面，基于Flask和原生JavaScript搭建了Web Console图形化管理界面。实验基于20个真实CVE漏洞和3种大语言模型进行300次独立测试，结果表明改进版总体成功率达到63.67\%，较原版提升约39个百分点，单次测试成本降低69\%，平均耗时缩短42\%。
}

% 毕业设计：图纸数量与正文字数（成果类型为 design 或 engineering 时填写，两者须同时填写）
% \infodrawings{5}
% \infowordcount{15000}

% 毕业论文：正文总字数（成果类型为 thesis 时填写）
\infothesiswords{约21000}

% 随附资料（每条用 \item 开头；不填则显示空白行）
\infomaterials{
  \item 毕业设计（论文）全套电子版；
  \item 程序源代码
  \item 附录
}
